% Options for packages loaded elsewhere
\PassOptionsToPackage{unicode}{hyperref}
\PassOptionsToPackage{hyphens}{url}
\documentclass[
  11pt,
  a4paper,
]{article}
\usepackage{xcolor}
\usepackage[margin=18mm]{geometry}
\usepackage{amsmath,amssymb}
\setcounter{secnumdepth}{-\maxdimen} % remove section numbering
\usepackage{iftex}
\ifPDFTeX
  \usepackage[T1]{fontenc}
  \usepackage[utf8]{inputenc}
  \usepackage{textcomp} % provide euro and other symbols
\else % if luatex or xetex
  \usepackage{unicode-math} % this also loads fontspec
  \defaultfontfeatures{Scale=MatchLowercase}
  \defaultfontfeatures[\rmfamily]{Ligatures=TeX,Scale=1}
\fi
\ifPDFTeX\else
  % xetex/luatex font selection
    \setmainfont[]{Noto Serif}
    \setsansfont[]{Noto Sans}
    \setmonofont[]{Noto Sans Mono}
  \setmathfont[]{Noto Sans Math}
\fi
% Use upquote if available, for straight quotes in verbatim environments
\IfFileExists{upquote.sty}{\usepackage{upquote}}{}
\IfFileExists{microtype.sty}{% use microtype if available
  \usepackage[]{microtype}
  \UseMicrotypeSet[protrusion]{basicmath} % disable protrusion for tt fonts
}{}
\makeatletter
\@ifundefined{KOMAClassName}{% if non-KOMA class
  \IfFileExists{parskip.sty}{%
    \usepackage{parskip}
  }{% else
    \setlength{\parindent}{0pt}
    \setlength{\parskip}{6pt plus 2pt minus 1pt}}
}{% if KOMA class
  \KOMAoptions{parskip=half}}
\makeatother
\ifLuaTeX
\usepackage[bidi=basic]{babel}
\else
\usepackage[bidi=default]{babel}
\fi
\babelprovide[main,import]{russian}
\ifPDFTeX
\else
\babelfont{rm}[]{Noto Serif}
\fi
% get rid of language-specific shorthands (see #6817):
\let\LanguageShortHands\languageshorthands
\def\languageshorthands#1{}
\setlength{\emergencystretch}{3em} % prevent overfull lines
\providecommand{\tightlist}{%
  \setlength{\itemsep}{0pt}\setlength{\parskip}{0pt}}
\usepackage{bookmark}
\IfFileExists{xurl.sty}{\usepackage{xurl}}{} % add URL line breaks if available
\urlstyle{same}
\hypersetup{
  pdftitle={Билеты по дискретным случайным процессам},
  pdflang={ru-RU},
  hidelinks,
  pdfcreator={LaTeX via pandoc}}

\title{Билеты по дискретным случайным процессам}
\author{}
\date{}

\begin{document}
\maketitle

\section{Билет 13. Связи математического ожидания и матрицы корреляции с
характеристической
функцией}\label{ux431ux438ux43bux435ux442-13.-ux441ux432ux44fux437ux438-ux43cux430ux442ux435ux43cux430ux442ux438ux447ux435ux441ux43aux43eux433ux43e-ux43eux436ux438ux434ux430ux43dux438ux44f-ux438-ux43cux430ux442ux440ux438ux446ux44b-ux43aux43eux440ux440ux435ux43bux44fux446ux438ux438-ux441-ux445ux430ux440ux430ux43aux442ux435ux440ux438ux441ux442ux438ux447ux435ux441ux43aux43eux439-ux444ux443ux43dux43aux446ux438ux435ux439}

Источники: Ширяев 1, гл. II, §12; лекции ДСП.

\subsection{1. Короткий
ответ}\label{ux43aux43eux440ux43eux442ux43aux438ux439-ux43eux442ux432ux435ux442}

Характеристическая функция конечномерного распределения кодирует не
только само распределение, но и моменты, если они существуют.

Для случайной величины \(X\):

\[
\varphi_X(t)=\mathbb E e^{itX}.
\]

Если \(\mathbb E|X|^k<\infty\), то \(\varphi_X\) имеет производные до
порядка \(k\), и

\[
\varphi_X^{(r)}(0)=i^r\mathbb E X^r,
\qquad r=1,\ldots,k.
\]

В частности,

\[
\mathbb EX=\frac{\varphi_X'(0)}{i}=-i\varphi_X'(0),
\]

а если существует второй момент, то

\[
\mathbb EX^2=-\varphi_X''(0),
\qquad
\operatorname{Var}X=-\varphi_X''(0)-(\mathbb EX)^2.
\]

Для случайного вектора \(X=(X_1,\ldots,X_n)^T\):

\[
\varphi_X(t)=\mathbb E e^{i\langle t,X\rangle},
\qquad t\in\mathbb R^n.
\]

Если \(\mathbb E|X_j|<\infty\), то

\[
\frac{\partial\varphi_X}{\partial t_j}(0)=i\mathbb EX_j.
\]

Если существуют вторые моменты, то

\[
\frac{\partial^2\varphi_X}{\partial t_i\partial t_j}(0)
=
-\mathbb E[X_iX_j].
\]

Значит средний вектор и матрица вторых моментов восстанавливаются из
производных характеристической функции в нуле:

\[
m_j=\mathbb EX_j=-i\partial_j\varphi_X(0),
\]

\[
R_{ij}=\mathbb E[X_iX_j]=-\partial_{ij}\varphi_X(0).
\]

Ковариационная матрица:

\[
\Sigma_{ij}
=
\operatorname{Cov}(X_i,X_j)
=
R_{ij}-m_im_j.
\]

Терминологически в курсах случайных процессов матрицу
\(R=(\mathbb E X_iX_j)\) часто называют корреляционной матрицей, а
\(\Sigma\) - ковариационной. Если величины центрированы, то
\(R=\Sigma\).

\subsection{2. Полный
ответ}\label{ux43fux43eux43bux43dux44bux439-ux43eux442ux432ux435ux442}

\subsubsection{2.1. Введение и связь с предыдущими
билетами}\label{ux432ux432ux435ux434ux435ux43dux438ux435-ux438-ux441ux432ux44fux437ux44c-ux441-ux43fux440ux435ux434ux44bux434ux443ux449ux438ux43cux438-ux431ux438ux43bux435ux442ux430ux43cux438}

Этот билет продолжает Билет 08 о характеристических функциях и Билет 11
о моментах через интегралы по распределению. В билете 8
характеристическая функция была введена как преобразование Фурье
конечномерного распределения:

\[
\varphi_X(t)
=
\mathbb E e^{i\langle t,X\rangle}
=
\int_{\mathbb R^n} e^{i\langle t,x\rangle}\,P_X(dx).
\]

Здесь нужно объяснить, как из этой функции получить математическое
ожидание, вторые моменты и корреляционную или ковариационную матрицу.

\subsubsection{2.2. Одномерный случай: моменты случайной
величины}\label{ux43eux434ux43dux43eux43cux435ux440ux43dux44bux439-ux441ux43bux443ux447ux430ux439-ux43cux43eux43cux435ux43dux442ux44b-ux441ux43bux443ux447ux430ux439ux43dux43eux439-ux432ux435ux43bux438ux447ux438ux43dux44b}

Начнем с одномерного случая. Пусть \(X\) - вещественная случайная
величина. Ее характеристическая функция:

\[
\varphi(t)=\mathbb E e^{itX}.
\]

Она существует всегда, потому что

\[
|e^{itX}|=1.
\]

Но производные характеристической функции и моменты требуют
дополнительных условий. Если \(\mathbb E|X|<\infty\), то можно
дифференцировать под знаком математического ожидания:

\[
\varphi'(t)=\mathbb E\left[iXe^{itX}\right].
\]

При \(t=0\):

\[
\varphi'(0)=i\mathbb EX.
\]

Поэтому

\[
\mathbb EX=\frac{\varphi'(0)}{i}=-i\varphi'(0).
\]

Если \(\mathbb EX^2<\infty\), то

\[
\varphi''(t)=\mathbb E\left[(iX)^2e^{itX}\right]
=
-\mathbb E\left[X^2e^{itX}\right],
\]

и при \(t=0\):

\[
\varphi''(0)=-\mathbb EX^2.
\]

Значит

\[
\mathbb EX^2=-\varphi''(0).
\]

Дисперсия:

\[
\operatorname{Var}X
=
\mathbb EX^2-(\mathbb EX)^2
=
-\varphi''(0)-(-i\varphi'(0))^2.
\]

Можно записать и проще:

\[
\operatorname{Var}X=-\varphi''(0)-(\mathbb EX)^2.
\]

Общая одномерная формула такая. Если

\[
\mathbb E|X|^k<\infty,
\]

то для \(r=1,\ldots,k\):

\[
\varphi^{(r)}(t)
=
\mathbb E\left[(iX)^r e^{itX}\right],
\]

а в нуле:

\[
\varphi^{(r)}(0)
=
i^r\mathbb EX^r.
\]

Отсюда

\[
\mathbb EX^r
=
\frac{\varphi^{(r)}(0)}{i^r}.
\]

\subsubsection{2.3. Случайный вектор: математическое ожидание и матрица
вторых
моментов}\label{ux441ux43bux443ux447ux430ux439ux43dux44bux439-ux432ux435ux43aux442ux43eux440-ux43cux430ux442ux435ux43cux430ux442ux438ux447ux435ux441ux43aux43eux435-ux43eux436ux438ux434ux430ux43dux438ux435-ux438-ux43cux430ux442ux440ux438ux446ux430-ux432ux442ux43eux440ux44bux445-ux43cux43eux43cux435ux43dux442ux43eux432}

Теперь случайный вектор. Пусть

\[
X=(X_1,\ldots,X_n)^T
\]

и \(P_X\) - его конечномерное распределение на \(\mathbb R^n\).
Характеристическая функция:

\[
\varphi(t)
=
\varphi_X(t_1,\ldots,t_n)
=
\mathbb E\exp\left(i\sum_{j=1}^n t_jX_j\right).
\]

Если компоненты интегрируемы, то частные производные первого порядка в
нуле дают математические ожидания:

\[
\partial_j\varphi(t)
=
\mathbb E\left[iX_j e^{i\langle t,X\rangle}\right],
\]

и

\[
\partial_j\varphi(0)=i\mathbb EX_j.
\]

Поэтому средний вектор

\[
m=\mathbb EX=(m_1,\ldots,m_n)^T
\]

получается по формуле

\[
m_j=-i\partial_j\varphi(0).
\]

Если существуют вторые моменты, то вторые частные производные дают
смешанные вторые моменты:

\[
\partial_{ij}\varphi(t)
=
\frac{\partial^2\varphi(t)}{\partial t_i\partial t_j}
=
\mathbb E\left[(iX_i)(iX_j)e^{i\langle t,X\rangle}\right]
=
-\mathbb E\left[X_iX_j e^{i\langle t,X\rangle}\right].
\]

При \(t=0\):

\[
\partial_{ij}\varphi(0)
=
-\mathbb E[X_iX_j].
\]

Обозначим матрицу вторых моментов:

\[
R=\mathbb E[XX^T],
\qquad
R_{ij}=\mathbb E[X_iX_j].
\]

Тогда

\[
R_{ij}=-\partial_{ij}\varphi(0).
\]

В терминах матрицы Гессе:

\[
R=-D^2\varphi(0),
\]

если \(D^2\varphi(0)\) понимается как матрица вторых частных
производных.

\subsubsection{2.4. Ковариационная матрица и
терминология}\label{ux43aux43eux432ux430ux440ux438ux430ux446ux438ux43eux43dux43dux430ux44f-ux43cux430ux442ux440ux438ux446ux430-ux438-ux442ux435ux440ux43cux438ux43dux43eux43bux43eux433ux438ux44f}

Ковариационная матрица:

\[
\Sigma
=
\operatorname{Cov}(X)
=
\mathbb E[(X-m)(X-m)^T].
\]

По формуле из Билета 11:

\[
\Sigma=R-mm^T.
\]

Следовательно, через характеристическую функцию:

\[
\Sigma_{ij}
=
-\partial_{ij}\varphi(0)-m_im_j.
\]

Если выразить \(m_i\) и \(m_j\) тоже через производные \(\varphi\), то

\[
\Sigma_{ij}
=
-\partial_{ij}\varphi(0)
+\partial_i\varphi(0)\partial_j\varphi(0).
\]

Проверка знака: так как \(\partial_i\varphi(0)=im_i\), то

\[
\partial_i\varphi(0)\partial_j\varphi(0)
=
-m_im_j.
\]

Поэтому последняя формула действительно равна

\[
-\partial_{ij}\varphi(0)-m_im_j.
\]

Отдельно полезно сказать про термин ``матрица корреляции''. В разных
курсах используются три близких объекта:

\begin{enumerate}
\def\labelenumi{\arabic{enumi}.}
\tightlist
\item
  матрица вторых моментов, или корреляционная матрица в теории случайных
  процессов:
\end{enumerate}

\[
R_{ij}=\mathbb E[X_iX_j];
\]

\begin{enumerate}
\def\labelenumi{\arabic{enumi}.}
\setcounter{enumi}{1}
\tightlist
\item
  ковариационная матрица:
\end{enumerate}

\[
\Sigma_{ij}=\operatorname{Cov}(X_i,X_j);
\]

\begin{enumerate}
\def\labelenumi{\arabic{enumi}.}
\setcounter{enumi}{2}
\tightlist
\item
  нормированная матрица коэффициентов корреляции:
\end{enumerate}

\[
\rho_{ij}
=
\frac{\operatorname{Cov}(X_i,X_j)}
{\sqrt{\operatorname{Var}X_i\,\operatorname{Var}X_j}},
\]

если дисперсии положительны.

В билете 13 обычно нужна связь характеристической функции именно с
математическим ожиданием и матрицей вторых моментов \(R\), а затем с
ковариационной матрицей \(\Sigma\). Если \(X\) центрирован, то

\[
m=0,
\qquad
\Sigma=R,
\]

и корреляционная матрица в смысле вторых моментов совпадает с
ковариационной.

\subsubsection{2.5. Приложение к конечномерным распределениям случайных
процессов}\label{ux43fux440ux438ux43bux43eux436ux435ux43dux438ux435-ux43a-ux43aux43eux43dux435ux447ux43dux43eux43cux435ux440ux43dux44bux43c-ux440ux430ux441ux43fux440ux435ux434ux435ux43bux435ux43dux438ux44fux43c-ux441ux43bux443ux447ux430ux439ux43dux44bux445-ux43fux440ux43eux446ux435ux441ux441ux43eux432}

Для конечномерного распределения случайного процесса запись такая же.
Если выбран набор моментов времени

\[
t_1,\ldots,t_n,
\]

и

\[
X=(X_{t_1},\ldots,X_{t_n})^T,
\]

то его конечномерная характеристическая функция:

\[
\varphi_{t_1,\ldots,t_n}(u_1,\ldots,u_n)
=
\mathbb E\exp\left(i\sum_{j=1}^n u_jX_{t_j}\right).
\]

Тогда

\[
\mathbb EX_{t_j}
=
-i\frac{\partial}{\partial u_j}
\varphi_{t_1,\ldots,t_n}(0),
\]

а корреляционная функция на выбранном конечном наборе индексов:

\[
R(t_i,t_j)
=
\mathbb E[X_{t_i}X_{t_j}]
=
-\frac{\partial^2}{\partial u_i\partial u_j}
\varphi_{t_1,\ldots,t_n}(0).
\]

Ковариационная функция:

\[
K(t_i,t_j)
=
\operatorname{Cov}(X_{t_i},X_{t_j})
=
R(t_i,t_j)-m(t_i)m(t_j).
\]

В стационарных задачах второго порядка это переходит в билеты Билет 17,
Билет 24 и Билет 25: среднее постоянно, а корреляция или ковариация
зависит от лага.

\subsubsection{2.6. Логарифм характеристической функции и
кумулянты}\label{ux43bux43eux433ux430ux440ux438ux444ux43c-ux445ux430ux440ux430ux43aux442ux435ux440ux438ux441ux442ux438ux447ux435ux441ux43aux43eux439-ux444ux443ux43dux43aux446ux438ux438-ux438-ux43aux443ux43cux443ux43bux44fux43dux442ux44b}

Иногда удобнее использовать логарифм характеристической функции. Так как
\(\varphi(0)=1\), в некоторой окрестности нуля при достаточной гладкости
можно рассматривать

\[
\psi(t)=\log\varphi(t).
\]

Тогда первые производные \(\psi\) дают среднее, а вторые - ковариации:

\[
\partial_j\psi(0)=i m_j,
\]

\[
\partial_{ij}\psi(0)=-\Sigma_{ij}.
\]

Это форма через кумулянты: первая кумулянта - среднее, вторая кумулянта
- ковариация. Для экзамена достаточно знать основную формулу через
производные \(\varphi\); логарифм можно добавить как сильное замечание.

\subsubsection{2.7. Главная идея
билета}\label{ux433ux43bux430ux432ux43dux430ux44f-ux438ux434ux435ux44f-ux431ux438ux43bux435ux442ux430}

Главная идея всего билета: характеристическая функция является
преобразованием Фурье распределения. Дифференцирование по параметрам
\(t_j\) выносит из экспоненты множители \(iX_j\). Поэтому производные в
нуле превращаются в моменты.

\subsection{3. Основные
определения}\label{ux43eux441ux43dux43eux432ux43dux44bux435-ux43eux43fux440ux435ux434ux435ux43bux435ux43dux438ux44f}

\subsubsection{Характеристическая функция случайной
величины}\label{ux445ux430ux440ux430ux43aux442ux435ux440ux438ux441ux442ux438ux447ux435ux441ux43aux430ux44f-ux444ux443ux43dux43aux446ux438ux44f-ux441ux43bux443ux447ux430ux439ux43dux43eux439-ux432ux435ux43bux438ux447ux438ux43dux44b}

\[
\varphi_X(t)=\mathbb E e^{itX},
\qquad t\in\mathbb R.
\]

\subsubsection{Характеристическая функция случайного
вектора}\label{ux445ux430ux440ux430ux43aux442ux435ux440ux438ux441ux442ux438ux447ux435ux441ux43aux430ux44f-ux444ux443ux43dux43aux446ux438ux44f-ux441ux43bux443ux447ux430ux439ux43dux43eux433ux43e-ux432ux435ux43aux442ux43eux440ux430}

\[
\varphi_X(t)=\mathbb E e^{i\langle t,X\rangle},
\qquad t\in\mathbb R^n.
\]

\subsubsection{Характеристическая функция конечномерного распределения
процесса}\label{ux445ux430ux440ux430ux43aux442ux435ux440ux438ux441ux442ux438ux447ux435ux441ux43aux430ux44f-ux444ux443ux43dux43aux446ux438ux44f-ux43aux43eux43dux435ux447ux43dux43eux43cux435ux440ux43dux43eux433ux43e-ux440ux430ux441ux43fux440ux435ux434ux435ux43bux435ux43dux438ux44f-ux43fux440ux43eux446ux435ux441ux441ux430}

Для процесса \((X_s)_{s\in T}\) и индексов \(s_1,\ldots,s_n\):

\[
\varphi_{s_1,\ldots,s_n}(u)
=
\mathbb E\exp\left(i\sum_{j=1}^n u_jX_{s_j}\right).
\]

\subsubsection{Математическое ожидание
вектора}\label{ux43cux430ux442ux435ux43cux430ux442ux438ux447ux435ux441ux43aux43eux435-ux43eux436ux438ux434ux430ux43dux438ux435-ux432ux435ux43aux442ux43eux440ux430}

\[
m=\mathbb EX,
\qquad
m_j=\mathbb EX_j.
\]

Оно существует, если все компоненты \(X_j\) интегрируемы.

\subsubsection{Матрица вторых
моментов}\label{ux43cux430ux442ux440ux438ux446ux430-ux432ux442ux43eux440ux44bux445-ux43cux43eux43cux435ux43dux442ux43eux432}

\[
R=\mathbb E[XX^T],
\qquad
R_{ij}=\mathbb E[X_iX_j].
\]

В теории случайных процессов ее часто называют корреляционной матрицей.

\subsubsection{Ковариационная
матрица}\label{ux43aux43eux432ux430ux440ux438ux430ux446ux438ux43eux43dux43dux430ux44f-ux43cux430ux442ux440ux438ux446ux430}

\[
\Sigma=\operatorname{Cov}(X)=\mathbb E[(X-m)(X-m)^T],
\]

то есть

\[
\Sigma_{ij}
=
\operatorname{Cov}(X_i,X_j)
=
\mathbb E[(X_i-m_i)(X_j-m_j)].
\]

\subsubsection{\texorpdfstring{Связь матриц \(R\) и
\(\Sigma\)}{Связь матриц R и \textbackslash Sigma}}\label{ux441ux432ux44fux437ux44c-ux43cux430ux442ux440ux438ux446-r-ux438-sigma}

\[
\Sigma=R-mm^T.
\]

Если \(m=0\), то

\[
\Sigma=R.
\]

\subsubsection{Нормированная матрица коэффициентов
корреляции}\label{ux43dux43eux440ux43cux438ux440ux43eux432ux430ux43dux43dux430ux44f-ux43cux430ux442ux440ux438ux446ux430-ux43aux43eux44dux444ux444ux438ux446ux438ux435ux43dux442ux43eux432-ux43aux43eux440ux440ux435ux43bux44fux446ux438ux438}

Если \(\operatorname{Var}X_i>0\), то

\[
\rho_{ij}
=
\frac{\Sigma_{ij}}
{\sqrt{\Sigma_{ii}\Sigma_{jj}}}.
\]

Диагональные элементы равны \(1\), а \(|\rho_{ij}|\le1\).

\subsubsection{Мультииндексная форма
момента}\label{ux43cux443ux43bux44cux442ux438ux438ux43dux434ux435ux43aux441ux43dux430ux44f-ux444ux43eux440ux43cux430-ux43cux43eux43cux435ux43dux442ux430}

Для \(\alpha=(\alpha_1,\ldots,\alpha_n)\):

\[
|\alpha|=\alpha_1+\cdots+\alpha_n,
\]

\[
\partial^\alpha
=
\frac{\partial^{|\alpha|}}
{\partial t_1^{\alpha_1}\cdots\partial t_n^{\alpha_n}},
\]

и при существовании соответствующего абсолютного момента:

\[
\partial^\alpha\varphi(0)
=
i^{|\alpha|}
\mathbb E\left[X_1^{\alpha_1}\cdots X_n^{\alpha_n}\right].
\]

\subsection{4. Основные теоремы и
свойства}\label{ux43eux441ux43dux43eux432ux43dux44bux435-ux442ux435ux43eux440ux435ux43cux44b-ux438-ux441ux432ux43eux439ux441ux442ux432ux430}

\subsubsection{Теорема 1. Моменты одномерной величины через производные
характеристической
функции}\label{ux442ux435ux43eux440ux435ux43cux430-1.-ux43cux43eux43cux435ux43dux442ux44b-ux43eux434ux43dux43eux43cux435ux440ux43dux43eux439-ux432ux435ux43bux438ux447ux438ux43dux44b-ux447ux435ux440ux435ux437-ux43fux440ux43eux438ux437ux432ux43eux434ux43dux44bux435-ux445ux430ux440ux430ux43aux442ux435ux440ux438ux441ux442ux438ux447ux435ux441ux43aux43eux439-ux444ux443ux43dux43aux446ux438ux438}

Если

\[
\mathbb E|X|^k<\infty,
\]

то \(\varphi_X\) имеет непрерывные производные до порядка \(k\), и

\[
\varphi_X^{(r)}(t)
=
\mathbb E[(iX)^r e^{itX}],
\qquad r=1,\ldots,k.
\]

В частности,

\[
\varphi_X^{(r)}(0)=i^r\mathbb EX^r.
\]

\subsubsection{Следствие. Среднее и
дисперсия}\label{ux441ux43bux435ux434ux441ux442ux432ux438ux435.-ux441ux440ux435ux434ux43dux435ux435-ux438-ux434ux438ux441ux43fux435ux440ux441ux438ux44f}

Если \(\mathbb E|X|<\infty\), то

\[
\mathbb EX=-i\varphi_X'(0).
\]

Если \(\mathbb EX^2<\infty\), то

\[
\mathbb EX^2=-\varphi_X''(0),
\]

\[
\operatorname{Var}X=-\varphi_X''(0)-(-i\varphi_X'(0))^2.
\]

\subsubsection{Теорема 2. Моменты вектора через частные
производные}\label{ux442ux435ux43eux440ux435ux43cux430-2.-ux43cux43eux43cux435ux43dux442ux44b-ux432ux435ux43aux442ux43eux440ux430-ux447ux435ux440ux435ux437-ux447ux430ux441ux442ux43dux44bux435-ux43fux440ux43eux438ux437ux432ux43eux434ux43dux44bux435}

Пусть \(X=(X_1,\ldots,X_n)^T\) и существует момент

\[
\mathbb E|X_1|^{\alpha_1}\cdots |X_n|^{\alpha_n}<\infty.
\]

Тогда

\[
\partial^\alpha\varphi_X(0)
=
i^{|\alpha|}
\mathbb E\left[X_1^{\alpha_1}\cdots X_n^{\alpha_n}\right].
\]

Для первого и второго порядков:

\[
\partial_j\varphi(0)=i\mathbb EX_j,
\]

\[
\partial_{ij}\varphi(0)=-\mathbb E[X_iX_j].
\]

\subsubsection{Свойство 3. Матрица вторых моментов через матрицу Гессе
характеристической
функции}\label{ux441ux432ux43eux439ux441ux442ux432ux43e-3.-ux43cux430ux442ux440ux438ux446ux430-ux432ux442ux43eux440ux44bux445-ux43cux43eux43cux435ux43dux442ux43eux432-ux447ux435ux440ux435ux437-ux43cux430ux442ux440ux438ux446ux443-ux433ux435ux441ux441ux435-ux445ux430ux440ux430ux43aux442ux435ux440ux438ux441ux442ux438ux447ux435ux441ux43aux43eux439-ux444ux443ux43dux43aux446ux438ux438}

Если все вторые моменты существуют, то

\[
R=-D^2\varphi(0),
\]

то есть

\[
R_{ij}=-\partial_{ij}\varphi(0).
\]

\subsubsection{Свойство 4. Ковариационная матрица через производные
характеристической
функции}\label{ux441ux432ux43eux439ux441ux442ux432ux43e-4.-ux43aux43eux432ux430ux440ux438ux430ux446ux438ux43eux43dux43dux430ux44f-ux43cux430ux442ux440ux438ux446ux430-ux447ux435ux440ux435ux437-ux43fux440ux43eux438ux437ux432ux43eux434ux43dux44bux435-ux445ux430ux440ux430ux43aux442ux435ux440ux438ux441ux442ux438ux447ux435ux441ux43aux43eux439-ux444ux443ux43dux43aux446ux438ux438}

Если все вторые моменты существуют, то

\[
\Sigma_{ij}
=
-\partial_{ij}\varphi(0)
+\partial_i\varphi(0)\partial_j\varphi(0).
\]

Эквивалентно:

\[
\Sigma=-D^2\varphi(0)-mm^T,
\qquad
m_j=-i\partial_j\varphi(0).
\]

\subsubsection{Свойство 5. Через логарифм характеристической
функции}\label{ux441ux432ux43eux439ux441ux442ux432ux43e-5.-ux447ux435ux440ux435ux437-ux43bux43eux433ux430ux440ux438ux444ux43c-ux445ux430ux440ux430ux43aux442ux435ux440ux438ux441ux442ux438ux447ux435ux441ux43aux43eux439-ux444ux443ux43dux43aux446ux438ux438}

Если \(\psi(t)=\log\varphi(t)\) определена и дважды дифференцируема в
окрестности нуля, то

\[
\partial_j\psi(0)=im_j,
\]

\[
\partial_{ij}\psi(0)=-\Sigma_{ij}.
\]

\subsubsection{Свойство 6. Неотрицательная определенность матриц второго
порядка}\label{ux441ux432ux43eux439ux441ux442ux432ux43e-6.-ux43dux435ux43eux442ux440ux438ux446ux430ux442ux435ux43bux44cux43dux430ux44f-ux43eux43fux440ux435ux434ux435ux43bux435ux43dux43dux43eux441ux442ux44c-ux43cux430ux442ux440ux438ux446-ux432ux442ux43eux440ux43eux433ux43e-ux43fux43eux440ux44fux434ux43aux430}

Матрица вторых моментов \(R\) неотрицательно определена:

\[
a^TRa
=
\mathbb E(a^TX)^2
\ge0.
\]

Ковариационная матрица \(\Sigma\) тоже неотрицательно определена:

\[
a^T\Sigma a
=
\operatorname{Var}(a^TX)
\ge0.
\]

Это согласуется с Билетом 12: \(R\) и \(\Sigma\) являются матрицами
Грама соответствующих случайных величин в \(L^2\).

\subsubsection{Свойство 7. Независимые
координаты}\label{ux441ux432ux43eux439ux441ux442ux432ux43e-7.-ux43dux435ux437ux430ux432ux438ux441ux438ux43cux44bux435-ux43aux43eux43eux440ux434ux438ux43dux430ux442ux44b}

Если \(X_1,\ldots,X_n\) независимы, то

\[
\varphi_X(t_1,\ldots,t_n)
=
\prod_{j=1}^n\varphi_{X_j}(t_j).
\]

Тогда смешанные моменты, если они существуют, раскладываются. Например:

\[
\mathbb E[X_iX_j]=\mathbb EX_i\,\mathbb EX_j,
\qquad i\ne j.
\]

Поэтому

\[
\operatorname{Cov}(X_i,X_j)=0
\]

для \(i\ne j\), если вторые моменты конечны.

\subsubsection{Свойство 8. Для гауссовского
вектора}\label{ux441ux432ux43eux439ux441ux442ux432ux43e-8.-ux434ux43bux44f-ux433ux430ux443ux441ux441ux43eux432ux441ux43aux43eux433ux43e-ux432ux435ux43aux442ux43eux440ux430}

Если

\[
X\sim N(m,\Sigma),
\]

то

\[
\varphi_X(t)
=
\exp\left(i\langle t,m\rangle-\frac12 t^T\Sigma t\right).
\]

Дифференцирование в нуле дает тот же средний вектор \(m\) и
ковариационную матрицу \(\Sigma\). Это связь с Билетом 18.

\subsection{5. Примеры}\label{ux43fux440ux438ux43cux435ux440ux44b}

\subsubsection{Пример 1. Вырожденная случайная
величина}\label{ux43fux440ux438ux43cux435ux440-1.-ux432ux44bux440ux43eux436ux434ux435ux43dux43dux430ux44f-ux441ux43bux443ux447ux430ux439ux43dux430ux44f-ux432ux435ux43bux438ux447ux438ux43dux430}

Пусть \(X=c\) почти наверное. Тогда

\[
\varphi(t)=e^{itc}.
\]

Производная:

\[
\varphi'(t)=ice^{itc},
\qquad
\varphi'(0)=ic.
\]

Значит

\[
\mathbb EX=-i\varphi'(0)=c.
\]

Вторая производная:

\[
\varphi''(0)=-c^2,
\]

поэтому

\[
\mathbb EX^2=-\varphi''(0)=c^2,
\]

и

\[
\operatorname{Var}X=c^2-c^2=0.
\]

\subsubsection{Пример 2. Бернуллиевская случайная
величина}\label{ux43fux440ux438ux43cux435ux440-2.-ux431ux435ux440ux43dux443ux43bux43bux438ux435ux432ux441ux43aux430ux44f-ux441ux43bux443ux447ux430ux439ux43dux430ux44f-ux432ux435ux43bux438ux447ux438ux43dux430}

Пусть

\[
P\{X=1\}=p,
\qquad
P\{X=0\}=q,
\qquad
p+q=1.
\]

Тогда

\[
\varphi(t)=q+pe^{it}.
\]

Первая производная:

\[
\varphi'(t)=ipe^{it},
\qquad
\varphi'(0)=ip.
\]

Значит

\[
\mathbb EX=p.
\]

Вторая производная:

\[
\varphi''(t)=-pe^{it},
\qquad
\varphi''(0)=-p.
\]

Поэтому

\[
\mathbb EX^2=p.
\]

Дисперсия:

\[
\operatorname{Var}X=p-p^2=pq.
\]

\subsubsection{Пример 3. Нормальная случайная
величина}\label{ux43fux440ux438ux43cux435ux440-3.-ux43dux43eux440ux43cux430ux43bux44cux43dux430ux44f-ux441ux43bux443ux447ux430ux439ux43dux430ux44f-ux432ux435ux43bux438ux447ux438ux43dux430}

Если

\[
X\sim N(m,\sigma^2),
\]

то

\[
\varphi(t)=\exp\left(imt-\frac12\sigma^2t^2\right).
\]

Логарифм:

\[
\psi(t)=\log\varphi(t)=imt-\frac12\sigma^2t^2.
\]

Тогда

\[
\psi'(0)=im,
\qquad
\psi''(0)=-\sigma^2.
\]

Значит

\[
\mathbb EX=m,
\qquad
\operatorname{Var}X=\sigma^2.
\]

\subsubsection{Пример 4. Дискретный двумерный
вектор}\label{ux43fux440ux438ux43cux435ux440-4.-ux434ux438ux441ux43aux440ux435ux442ux43dux44bux439-ux434ux432ux443ux43cux435ux440ux43dux44bux439-ux432ux435ux43aux442ux43eux440}

Пусть \(X=(X_1,X_2)\) принимает значения

\[
(0,0),\ (1,0),\ (0,1)
\]

с вероятностями \(p_0,p_1,p_2\), где \(p_0+p_1+p_2=1\). Тогда

\[
\varphi(t_1,t_2)
=
p_0+p_1e^{it_1}+p_2e^{it_2}.
\]

Средние:

\[
\partial_1\varphi(0)=ip_1,
\qquad
\partial_2\varphi(0)=ip_2,
\]

поэтому

\[
\mathbb EX_1=p_1,
\qquad
\mathbb EX_2=p_2.
\]

Смешанный второй момент:

\[
\partial_{12}\varphi(0)=0,
\]

значит

\[
\mathbb E[X_1X_2]=0.
\]

Ковариация:

\[
\operatorname{Cov}(X_1,X_2)
=
0-p_1p_2
=
-p_1p_2.
\]

Это хороший пример того, что смешанный второй момент и ковариация - не
одно и то же.

\subsubsection{Пример 5. Независимые
координаты}\label{ux43fux440ux438ux43cux435ux440-5.-ux43dux435ux437ux430ux432ux438ux441ux438ux43cux44bux435-ux43aux43eux43eux440ux434ux438ux43dux430ux442ux44b}

Пусть \(X_1\) и \(X_2\) независимы. Тогда

\[
\varphi_{(X_1,X_2)}(t_1,t_2)
=
\varphi_{X_1}(t_1)\varphi_{X_2}(t_2).
\]

Если существуют вторые моменты, то

\[
\mathbb E[X_1X_2]
=
\mathbb EX_1\,\mathbb EX_2,
\]

и поэтому

\[
\operatorname{Cov}(X_1,X_2)=0.
\]

Но обратное неверно: нулевая ковариация не означает независимость.

\subsubsection{Пример 6. Конечномерное распределение
процесса}\label{ux43fux440ux438ux43cux435ux440-6.-ux43aux43eux43dux435ux447ux43dux43eux43cux435ux440ux43dux43eux435-ux440ux430ux441ux43fux440ux435ux434ux435ux43bux435ux43dux438ux435-ux43fux440ux43eux446ux435ux441ux441ux430}

Пусть есть процесс \((X_t)\) и выбраны \(s,t\). Тогда

\[
\varphi_{s,t}(u,v)
=
\mathbb E e^{i(uX_s+vX_t)}.
\]

Если существуют вторые моменты, то

\[
\mathbb EX_s=-i\partial_u\varphi_{s,t}(0,0),
\]

\[
\mathbb E[X_sX_t]
=
-\partial_{uv}\varphi_{s,t}(0,0).
\]

Это и есть связь конечномерной характеристической функции с
корреляционной функцией процесса.

\subsection{6. Схемы доказательств /
обоснований}\label{ux441ux445ux435ux43cux44b-ux434ux43eux43aux430ux437ux430ux442ux435ux43bux44cux441ux442ux432-ux43eux431ux43eux441ux43dux43eux432ux430ux43dux438ux439}

\textbf{Почему производная дает математическое ожидание.}

Пусть \(\mathbb E|X|<\infty\). Для фиксированного \(t\) формально:

\[
\frac{d}{dt}e^{itX}=iXe^{itX}.
\]

Чтобы вынести производную под знак ожидания, используем доминирование.
Например, разностное отношение:

\[
\frac{e^{i(t+h)X}-e^{itX}}{h}
=
e^{itX}\frac{e^{ihX}-1}{h}.
\]

Поточечно при \(h\to0\) оно стремится к

\[
iXe^{itX}.
\]

Кроме того,

\[
\left|\frac{e^{ihX}-1}{h}\right|
\le |X|
\]

из оценки \(|e^{iy}-1|\le |y|\). Так как \(\mathbb E|X|<\infty\), можно
применить теорему Лебега о доминированной сходимости. Получаем

\[
\varphi'(t)=\mathbb E[iXe^{itX}].
\]

При \(t=0\):

\[
\varphi'(0)=i\mathbb EX.
\]

\textbf{Почему вторая производная дает второй момент.}

Если \(\mathbb EX^2<\infty\), то аналогично дифференцируем еще раз:

\[
\varphi''(t)
=
\mathbb E[(iX)^2e^{itX}]
=
-\mathbb E[X^2e^{itX}].
\]

При \(t=0\):

\[
\varphi''(0)=-\mathbb EX^2.
\]

Знак минус появляется из-за \(i^2=-1\).

\textbf{Доказательная схема для вектора.}

Для \(X=(X_1,\ldots,X_n)\):

\[
\varphi(t)=\mathbb E e^{i(t_1X_1+\cdots+t_nX_n)}.
\]

Дифференцирование по \(t_j\) выносит множитель \(iX_j\):

\[
\partial_j\varphi(t)
=
\mathbb E[iX_je^{i\langle t,X\rangle}].
\]

Дифференцирование по \(t_i\) и \(t_j\) выносит множитель

\[
(iX_i)(iX_j)=-X_iX_j.
\]

Отсюда:

\[
\partial_{ij}\varphi(0)=-\mathbb E[X_iX_j].
\]

\textbf{Почему ковариационная матрица неотрицательно определена.}

Для любого \(a\in\mathbb R^n\):

\[
a^T\Sigma a
=
\sum_{i,j}a_i a_j\operatorname{Cov}(X_i,X_j)
=
\operatorname{Var}\left(\sum_i a_iX_i\right)
\ge0.
\]

Это также следует из геометрии \(L^2\): ковариационная матрица - матрица
Грама центрированных величин.

\textbf{Почему через логарифм получается ковариация.}

Пусть

\[
\psi=\log\varphi.
\]

Тогда

\[
\partial_i\psi=\frac{\partial_i\varphi}{\varphi}.
\]

В нуле \(\varphi(0)=1\), поэтому

\[
\partial_i\psi(0)=\partial_i\varphi(0)=im_i.
\]

Для второй производной:

\[
\partial_{ij}\psi(0)
=
\partial_{ij}\varphi(0)
-\partial_i\varphi(0)\partial_j\varphi(0).
\]

Подставляем:

\[
\partial_{ij}\varphi(0)=-\mathbb E[X_iX_j],
\]

\[
\partial_i\varphi(0)\partial_j\varphi(0)
=
(im_i)(im_j)
=
-m_im_j.
\]

Получаем

\[
\partial_{ij}\psi(0)
=
-\mathbb E[X_iX_j]+m_im_j
=
-\operatorname{Cov}(X_i,X_j).
\]

\subsection{7. Что написать на
доске}\label{ux447ux442ux43e-ux43dux430ux43fux438ux441ux430ux442ux44c-ux43dux430-ux434ux43eux441ux43aux435}

\begin{enumerate}
\def\labelenumi{\arabic{enumi}.}
\tightlist
\item
  Характеристическая функция:
\end{enumerate}

\[
\varphi_X(t)=\mathbb E e^{i\langle t,X\rangle}.
\]

\begin{enumerate}
\def\labelenumi{\arabic{enumi}.}
\setcounter{enumi}{1}
\tightlist
\item
  Одномерные моменты:
\end{enumerate}

\[
\varphi_X^{(k)}(0)=i^k\mathbb EX^k.
\]

\begin{enumerate}
\def\labelenumi{\arabic{enumi}.}
\setcounter{enumi}{2}
\tightlist
\item
  Средний вектор:
\end{enumerate}

\[
m_j=\mathbb EX_j=-i\partial_j\varphi(0).
\]

\begin{enumerate}
\def\labelenumi{\arabic{enumi}.}
\setcounter{enumi}{3}
\tightlist
\item
  Матрица вторых моментов:
\end{enumerate}

\[
R_{ij}=\mathbb E[X_iX_j]=-\partial_{ij}\varphi(0).
\]

\begin{enumerate}
\def\labelenumi{\arabic{enumi}.}
\setcounter{enumi}{4}
\tightlist
\item
  Ковариационная матрица:
\end{enumerate}

\[
\Sigma_{ij}=R_{ij}-m_im_j.
\]

\begin{enumerate}
\def\labelenumi{\arabic{enumi}.}
\setcounter{enumi}{5}
\tightlist
\item
  Через производные \(\varphi\):
\end{enumerate}

\[
\Sigma_{ij}
=
-\partial_{ij}\varphi(0)
+\partial_i\varphi(0)\partial_j\varphi(0).
\]

\begin{enumerate}
\def\labelenumi{\arabic{enumi}.}
\setcounter{enumi}{6}
\tightlist
\item
  Через логарифм:
\end{enumerate}

\[
\partial_{ij}\log\varphi(0)=-\Sigma_{ij}.
\]

\begin{enumerate}
\def\labelenumi{\arabic{enumi}.}
\setcounter{enumi}{7}
\tightlist
\item
  Для процесса:
\end{enumerate}

\[
R(s,t)=\mathbb E[X_sX_t]
=
-\partial_{uv}\varphi_{s,t}(0,0).
\]

\subsection{8. Типичные дополнительные
вопросы}\label{ux442ux438ux43fux438ux447ux43dux44bux435-ux434ux43eux43fux43eux43bux43dux438ux442ux435ux43bux44cux43dux44bux435-ux432ux43eux43fux440ux43eux441ux44b}

\textbf{Вопрос. Характеристическая функция всегда существует?}

Да. Для любого \(t\):

\[
|e^{i\langle t,X\rangle}|=1,
\]

поэтому математическое ожидание существует как ожидание ограниченной
комплексной случайной величины.

\textbf{Вопрос. Если характеристическая функция существует всегда,
значит ли это, что все моменты существуют?}

Нет. Существование \(\varphi\) не означает существования моментов.
Например, у распределения Коши характеристическая функция существует, но
математическое ожидание не существует.

\textbf{Вопрос. Что нужно для формулы \(\varphi'(0)=i\mathbb EX\)?}

Достаточно \(\mathbb E|X|<\infty\). Тогда можно дифференцировать под
знаком ожидания.

\textbf{Вопрос. Что нужно для формулы \(\varphi''(0)=-\mathbb EX^2\)?}

Достаточно \(\mathbb EX^2<\infty\).

\textbf{Вопрос. Почему во второй производной стоит минус?}

Потому что

\[
\frac{d^2}{dt^2}e^{itX}
=(iX)^2e^{itX}
=-X^2e^{itX}.
\]

\textbf{Вопрос. Как получить математическое ожидание вектора?}

По компонентам:

\[
\mathbb EX_j=-i\partial_j\varphi(0).
\]

\textbf{Вопрос. Как получить смешанный момент \(\mathbb E[X_iX_j]\)?}

Через смешанную вторую производную:

\[
\mathbb E[X_iX_j]=-\partial_{ij}\varphi(0).
\]

\textbf{Вопрос. Чем отличается матрица вторых моментов от ковариационной
матрицы?}

Матрица вторых моментов:

\[
R_{ij}=\mathbb E[X_iX_j].
\]

Ковариационная матрица:

\[
\Sigma_{ij}=\mathbb E[(X_i-m_i)(X_j-m_j)].
\]

Связь:

\[
\Sigma=R-mm^T.
\]

\textbf{Вопрос. Когда корреляционная матрица совпадает с
ковариационной?}

Когда вектор центрирован:

\[
\mathbb EX=0.
\]

Тогда

\[
R=\Sigma.
\]

\textbf{Вопрос. Что дает логарифм характеристической функции?}

Он дает кумулянты. Первая производная в нуле дает среднее, а вторая -
ковариационную матрицу:

\[
\partial_j\log\varphi(0)=im_j,
\qquad
\partial_{ij}\log\varphi(0)=-\Sigma_{ij}.
\]

\textbf{Вопрос. Как это связано с гауссовским распределением?}

Для гауссовского вектора

\[
\varphi(t)=\exp\left(i\langle t,m\rangle-\frac12t^T\Sigma t\right).
\]

Из логарифма сразу видны \(m\) и \(\Sigma\).

\textbf{Вопрос. Как это связано со случайными процессами?}

Если взять конечномерное распределение

\[
(X_{t_1},\ldots,X_{t_n}),
\]

то производные его характеристической функции в нуле дают средние
\(m(t_i)\) и корреляции \(R(t_i,t_j)\).

\subsection{9. Типичные
ошибки}\label{ux442ux438ux43fux438ux447ux43dux44bux435-ux43eux448ux438ux431ux43aux438}

\begin{enumerate}
\def\labelenumi{\arabic{enumi}.}
\tightlist
\item
  Писать \(\varphi''(0)=\mathbb EX^2\) без минуса.
\item
  Забывать множитель \(i\) в первой производной.
\item
  Считать, что характеристическая функция существует только при
  существовании моментов.
\item
  Считать, что существование характеристической функции автоматически
  дает существование всех моментов.
\item
  Путать матрицу вторых моментов \(R=\mathbb E[XX^T]\) и ковариационную
  матрицу \(\Sigma=R-mm^T\).
\item
  Не центрировать величины при переходе от вторых моментов к
  ковариациям.
\item
  Путать корреляционную матрицу в смысле \(R_{ij}=\mathbb E[X_iX_j]\) с
  нормированной матрицей коэффициентов корреляции \(\rho_{ij}\).
\item
  Забывать, что формулы с производными в нуле требуют существования
  соответствующих моментов.
\item
  Дифференцировать по одному аргументу, когда нужен смешанный момент
  \(\mathbb E[X_iX_j]\).
\item
  Делать вывод о независимости из нулевой ковариации.
\end{enumerate}

\subsection{10. Связи с другими
билетами}\label{ux441ux432ux44fux437ux438-ux441-ux434ux440ux443ux433ux438ux43cux438-ux431ux438ux43bux435ux442ux430ux43cux438}

\textbf{Билет 08.} Там вводятся характеристические функции конечномерных
распределений и их основные свойства; здесь используются производные
этих функций.

\textbf{Билет 11.} Моменты, дисперсии и ковариации выражаются через
интегралы по распределению; характеристическая функция является другим
способом кодировать те же моменты.

\textbf{Билет 12.} Ковариационная матрица - матрица Грама центрированных
случайных величин в \(L^2\).

\textbf{Билет 14.} Коэффициенты корреляции получаются из ковариационной
матрицы нормировкой на стандартные отклонения.

\textbf{Билет 15.} Согласованные конечномерные распределения можно
задавать через согласованные конечномерные характеристические функции.

\textbf{Билет 18.} Гауссовское распределение полностью задается средним
и ковариационной матрицей, которые видны в характеристической функции.

\textbf{Билет 24.} Корреляционная функция процесса - это второй момент
или ковариация, получаемая из конечномерных характеристических функций.

\textbf{Билет 25.} Спектральные характеристики связаны с корреляционной
функцией как с коэффициентами Фурье.

\subsection{11. Минимум на
удовлетворительно}\label{ux43cux438ux43dux438ux43cux443ux43c-ux43dux430-ux443ux434ux43eux432ux43bux435ux442ux432ux43eux440ux438ux442ux435ux43bux44cux43dux43e}

Нужно уметь сказать:

\begin{enumerate}
\def\labelenumi{\arabic{enumi}.}
\tightlist
\item
  Характеристическая функция:
\end{enumerate}

\[
\varphi_X(t)=\mathbb E e^{i\langle t,X\rangle}.
\]

\begin{enumerate}
\def\labelenumi{\arabic{enumi}.}
\setcounter{enumi}{1}
\item
  Если существуют моменты, то производные в нуле дают эти моменты.
\item
  Для случайной величины:
\end{enumerate}

\[
\varphi'(0)=i\mathbb EX,
\qquad
\varphi''(0)=-\mathbb EX^2.
\]

\begin{enumerate}
\def\labelenumi{\arabic{enumi}.}
\setcounter{enumi}{3}
\tightlist
\item
  Для вектора:
\end{enumerate}

\[
\partial_j\varphi(0)=i\mathbb EX_j,
\qquad
\partial_{ij}\varphi(0)=-\mathbb E[X_iX_j].
\]

\begin{enumerate}
\def\labelenumi{\arabic{enumi}.}
\setcounter{enumi}{4}
\tightlist
\item
  Ковариационная матрица:
\end{enumerate}

\[
\Sigma_{ij}=\mathbb E[X_iX_j]-\mathbb EX_i\,\mathbb EX_j.
\]

Главная ловушка: во второй производной стоит минус.

\subsection{12. Минимум на
хорошо}\label{ux43cux438ux43dux438ux43cux443ux43c-ux43dux430-ux445ux43eux440ux43eux448ux43e}

Помимо минимума, нужно:

\begin{enumerate}
\def\labelenumi{\arabic{enumi}.}
\tightlist
\item
  назвать условие существования соответствующих моментов;
\item
  объяснить, почему дифференцирование выносит множитель \(iX_j\);
\item
  различить матрицу вторых моментов \(R\) и ковариационную матрицу
  \(\Sigma\);
\item
  записать формулы
\end{enumerate}

\[
m_j=-i\partial_j\varphi(0),
\]

\[
R_{ij}=-\partial_{ij}\varphi(0),
\]

\[
\Sigma=R-mm^T;
\]

\begin{enumerate}
\def\labelenumi{\arabic{enumi}.}
\setcounter{enumi}{4}
\tightlist
\item
  привести пример Бернулли или нормального распределения.
\end{enumerate}

\subsection{13. Что добавить для
отлично}\label{ux447ux442ux43e-ux434ux43eux431ux430ux432ux438ux442ux44c-ux434ux43bux44f-ux43eux442ux43bux438ux447ux43dux43e}

Для сильного ответа стоит добавить:

\begin{enumerate}
\def\labelenumi{\arabic{enumi}.}
\tightlist
\item
  мультииндексную формулу
\end{enumerate}

\[
\partial^\alpha\varphi(0)
=
i^{|\alpha|}
\mathbb E[X^\alpha];
\]

\begin{enumerate}
\def\labelenumi{\arabic{enumi}.}
\setcounter{enumi}{1}
\tightlist
\item
  доказательную схему через теорему Лебега о доминированной сходимости;
\item
  формулу через логарифм характеристической функции:
\end{enumerate}

\[
\partial_{ij}\log\varphi(0)=-\Sigma_{ij};
\]

\begin{enumerate}
\def\labelenumi{\arabic{enumi}.}
\setcounter{enumi}{3}
\tightlist
\item
  связь с конечномерными распределениями процесса:
\end{enumerate}

\[
R(s,t)=-\partial_{uv}\varphi_{s,t}(0,0);
\]

\begin{enumerate}
\def\labelenumi{\arabic{enumi}.}
\setcounter{enumi}{4}
\tightlist
\item
  неотрицательную определенность матриц \(R\) и \(\Sigma\);
\item
  связь с гауссовским вектором:
\end{enumerate}

\[
\varphi(t)=
\exp\left(i\langle t,m\rangle-\frac12t^T\Sigma t\right).
\]

\subsection{14. Устная версия ответа на 5
минут}\label{ux443ux441ux442ux43dux430ux44f-ux432ux435ux440ux441ux438ux44f-ux43eux442ux432ux435ux442ux430-ux43dux430-5-ux43cux438ux43dux443ux442}

Сначала я напоминаю определение характеристической функции. Для
случайного вектора

\[
X=(X_1,\ldots,X_n)
\]

она равна

\[
\varphi_X(t)=\mathbb E e^{i\langle t,X\rangle}.
\]

Она существует всегда, потому что экспонента имеет модуль \(1\). Но
производные в нуле связаны с моментами только при существовании
соответствующих моментов.

В одномерном случае, если \(\mathbb E|X|^k<\infty\), то

\[
\varphi_X^{(r)}(0)=i^r\mathbb EX^r,
\qquad r\le k.
\]

В частности,

\[
\varphi'(0)=i\mathbb EX,
\qquad
\varphi''(0)=-\mathbb EX^2.
\]

Значит

\[
\mathbb EX=-i\varphi'(0),
\qquad
\operatorname{Var}X=-\varphi''(0)-(\mathbb EX)^2.
\]

Для вектора дифференцируем по координатам:

\[
\partial_j\varphi(0)=i\mathbb EX_j,
\]

\[
\partial_{ij}\varphi(0)=-\mathbb E[X_iX_j].
\]

Поэтому средний вектор и матрица вторых моментов находятся как

\[
m_j=-i\partial_j\varphi(0),
\]

\[
R_{ij}=-\partial_{ij}\varphi(0).
\]

Дальше ковариационная матрица:

\[
\Sigma_{ij}
=
\operatorname{Cov}(X_i,X_j)
=
R_{ij}-m_im_j.
\]

Если вектор центрирован, то \(m=0\), и корреляционная матрица вторых
моментов совпадает с ковариационной.

Доказательная идея простая: при дифференцировании

\[
e^{i\langle t,X\rangle}
\]

по \(t_j\) появляется множитель \(iX_j\), а при втором дифференцировании
по \(t_i,t_j\) появляется \((iX_i)(iX_j)=-X_iX_j\). При существовании
нужных моментов дифференцирование под знаком ожидания оправдывается
доминированной сходимостью.

Завершаю примером. Для Бернулли

\[
\varphi(t)=q+pe^{it}.
\]

Тогда

\[
\varphi'(0)=ip,
\qquad
\varphi''(0)=-p,
\]

поэтому

\[
\mathbb EX=p,
\qquad
\operatorname{Var}X=p-p^2.
\]

Главные ошибки: потерять минус во второй производной, забыть множитель
\(i\), перепутать вторые моменты и ковариации.

\subsection{15. Если осталось 2 минуты перед
экзаменом}\label{ux435ux441ux43bux438-ux43eux441ux442ux430ux43bux43eux441ux44c-2-ux43cux438ux43dux443ux442ux44b-ux43fux435ux440ux435ux434-ux44dux43aux437ux430ux43cux435ux43dux43eux43c}

Запомнить четыре строки.

\textbf{Характеристическая функция:}

\[
\varphi_X(t)=\mathbb E e^{i\langle t,X\rangle}.
\]

\textbf{Среднее:}

\[
m_j=\mathbb EX_j=-i\partial_j\varphi(0).
\]

\textbf{Вторые моменты:}

\[
R_{ij}=\mathbb E[X_iX_j]=-\partial_{ij}\varphi(0).
\]

\textbf{Ковариации:}

\[
\Sigma_{ij}=R_{ij}-m_im_j.
\]

Главная ловушка: характеристическая функция существует всегда, но
производные-моменты требуют существования соответствующих моментов; во
второй производной стоит минус.

\newpage

\end{document}
